% This file was converted to LaTeX by Writer2LaTeX ver. 1.9.9
% see http://writer2latex.sourceforge.net for more info
\documentclass{article}
\usepackage{calc,amsmath,amssymb,amsfonts}
\usepackage[LGR,T1]{fontenc}
\usepackage[greek,italian]{babel}
\usepackage[style=numeric,backend=biber]{biblatex}
\author{HP}
\date{2026-08-19}
\begin{document}
Ci sono ricordi che il tempo sfuma lentamente, fino a renderli quasi irriconoscibili. Altri, invece, rimangono
sorprendentemente nitidi. Il mio primo incontro con un videogioco appartiene a questa seconda categoria.\newline
\textbf{Era il 1980.}\newline
Non saprei dire esattamente quale giorno fosse, ma se chiudo gli occhi riesco ancora a rivedere la stanza, immersa nella
luce soffusa che filtrava dalle tapparelle socchiuse per tenere fuori il caldo torrido dei pomeriggi estivi, e perfino
la sensazione provata in quell'istante.\newline
Frequentavo la scuola elementare quando un pomeriggio andai a trovare il mio amico Emanuele. Era una visita come tante
altre. Avremmo trascorso qualche ora insieme, giocando o semplicemente facendo quello che fanno due ragazzi della
stessa età, compresi quegli scherzi al telefono che facevano tanto arrabbiare i nostri parenti. Non avevo la minima
idea, invece, che quel pomeriggio avrebbe lasciato un segno destinato a rimanere con me per tutta la vita.\newline
Entrai nella sua cameretta. Ricordo una scrivania appoggiata contro la parete, un televisore acceso e lui che, con
naturalezza, prese in mano un joystick. Sullo schermo comparvero due barre bianche e un piccolo quadrato luminoso che
continuava a rimbalzare da una parte all'altra.\newline
\textbf{Era Pong.}\newline
Oggi potrebbe sembrare poco più di una curiosità storica. Ma allora era qualcosa di completamente diverso. Per me non
era soltanto un videogioco. Era una domanda.\newline
Guardavo quella piccola pallina cambiare direzione, rimbalzare, assegnare un punto, ricominciare la partita. Più
osservavo lo schermo, meno mi interessava giocare. Anni prima l$\text{\textgreek{’}}$avevo già visto nel bar del mio
paese ma a casa, dentro una scatola di plastica, mai! Continuavo a chiedermi come fosse possibile. Come faceva quel
televisore a sapere dove si trovasse la pallina? Come decideva quando invertire la direzione? Come riconosceva una
collisione? Come stabiliva il punteggio?\newline
Quando tornai a casa non avevo ancora trovato una risposta. Ma avevo trovato la domanda che avrebbe accompagnato gran
parte della mia vita.\newline
\textbf{Come fa?}\newline
Nelle settimane successive quella domanda non smise più di presentarsi. Ogni volta che osservavo un videogioco, ogni
volta che sentivo parlare di programmazione, tornava puntualmente a bussare nella mia mente.\newline
Non sapevo ancora che la risposta sarebbe arrivata pochi anni dopo, grazie a un computer che portava un numero
apparentemente insignificante.\newline
\textbf{Sessantaquattro.}\newline
Quel computer entrò nella mia vita quando ero ancora un ragazzo e, quasi senza accorgermene, cambiò il modo in cui
osservavo il mondo. Non mi limitavo più a utilizzare ciò che avevo davanti. Cominciai a chiedermi come fosse stato
costruito. Ogni comando imparato, ogni programma scritto, ogni errore corretto sembravano aggiungere un piccolo
tassello a una risposta che appariva ancora lontana.\newline
Come molti ragazzi della mia generazione, anch'io iniziai presto a coltivare un sogno. Non desideravo semplicemente
imparare il BASIC e non desideravo soltanto giocare. Volevo creare un videogioco vero, uno di quelli che riuscivano a
raccontare una storia, a dare vita a un personaggio con il movimento del joystick, a trasformare poche decine di
kilobyte in un mondo capace di far immaginare molto più di quanto mostrasse sullo schermo.\newline
Quel videogioco, però, non arrivò mai. Eppure avevamo provato perfino a immaginarlo. Con Emanuele, Vito e Gianni.
Pensavamo a un adventure, convinti che fosse più semplice da realizzare. Senza rendercene conto avevamo già formato una
piccola startup, molti anni prima che quella parola entrasse nel nostro vocabolario: io ed Emanuele avremmo
programmato, Vito avrebbe trasformato la sua passione per i fumetti in grafica e sceneggiatura, Gianni avrebbe pensato
alla musica e agli effetti sonori. Ne parlavamo molto. Ma non scrivemmo mai una sola riga di quel gioco.\newline
La vita, infatti, prese altre strade. Gli studi, il lavoro, le responsabilità e i progetti professionali di ciascuno
finirono, poco alla volta, per occupare tutto lo spazio. Il Commodore 64 venne riposto in una scatola insieme ai suoi
cavi, al Datassette, ai manuali e ai joystick consumati dall'uso.\newline
Pensavo che quella storia fosse finita.\newline
\textbf{Mi sbagliavo.}\newline
Molti anni dopo mi capitò di ritrovare quella scatola. Non ricordo nemmeno cosa stessi cercando. Forse una vecchia
valigia. Forse alcuni album fotografici. Ricordo soltanto la polvere, il cartone ormai ingiallito e la sensazione di
avere improvvisamente aperto una finestra sul passato.\newline
Dentro non c'era soltanto un vecchio computer, c'era una parte della mia storia. Presi il Commodore 64 tra le mani e
rimasi per qualche minuto in silenzio. Poi sorrisi.\newline
\textbf{$\text{\textgreek{«}}$Ci ho messo quarant'anni.$\text{\textgreek{»}}$}\newline
Quarant'anni per capire che esisteva ancora una promessa rimasta incompiuta.\newline
Il bambino che aveva desiderato con tutte le sue forze quella macchina aveva visto realizzarsi molto più di quanto
potesse immaginare. Quel computer gli aveva fatto scoprire la programmazione, lo aveva accompagnato nella scelta degli
studi e, con il tempo, aveva trasformato una curiosità in una professione. Aveva anche dato un senso al sacrificio dei
genitori che avevano deciso di affrontare uno sforzo economico importante per acquistarglielo, mantenendo la promessa,
fatta quasi senza rendersene conto, che quel computer non sarebbe stato soltanto un gioco.\newline
Ma un$\text{\textgreek{’}}$altra promessa era rimasta indietro: realizzare quel videogioco che aveva immaginato tante
volte e che non aveva mai avuto il coraggio, o forse semplicemente l'occasione, di cominciare.\newline
È da quel momento che nasce questo libro. Non dall'idea di scrivere un manuale sul BASIC del Commodore 64 e nemmeno dal
desiderio di raccogliere alcuni ricordi personali. Entrambe queste componenti sono presenti, ma nessuna rappresenta il
vero motivo per cui queste pagine esistono. Questo libro nasce dal desiderio di chiudere un cerchio rimasto aperto per
oltre quarant'anni.\newline
\textbf{Ma non voglio farlo da solo.}\newline
Vorrei che il lettore percorresse questo cammino insieme a me. Non importa se conserva ancora il proprio Commodore 64
perfettamente funzionante, se lo ha ritrovato dopo decenni in una cantina o se lo conoscerà soltanto attraverso un
emulatore. Non importa nemmeno se ha vissuto gli anni Ottanta o se appartiene a una generazione che li ha conosciuti
soltanto attraverso i racconti.\newline
Quello che conta è la curiosità. La stessa curiosità che nasce quando ci si trova davanti a qualcosa che sembra quasi
magico e si decide di non accontentarsi di usarlo, ma di capire come sia stato costruito.\newline
\textbf{Per questo motivo non affronteremo questo percorso come un manuale tradizionale.}\newline
Non sarà la teoria a guidare il progetto. Sarà il progetto stesso a indicarci, passo dopo passo, quali domande porci,
quali strumenti imparare e perché ci servono. Costruiremo un videogioco così come avrebbe potuto nascere negli anni
ottanta: un problema alla volta, una scelta alla volta, un errore alla volta. Ogni nuova conoscenza nascerà da una
necessità concreta. Ogni difficoltà diventerà l'occasione per imparare qualcosa di nuovo. Ogni soluzione farà crescere
il progetto.\newline
Il lettore non assisterà semplicemente allo sviluppo del programma, ma ne farà parte. Vedrà il videogioco nascere,
cambiare, sbagliare e migliorare. Scoprirà che i limiti del Commodore 64 non erano soltanto ostacoli, ma strumenti che
insegnavano a ragionare. Capirà che programmare significa molto più che impartire istruzioni a una macchina: significa
imparare a osservare un problema, formulare un'ipotesi e verificare se quell'idea sia davvero in grado di
funzionare.\newline
\textbf{Accanto al percorso tecnico ne scorrerà un altro, più personale.}\newline
È la storia delle persone che hanno alimentato quella curiosità, degli incontri che hanno lasciato un segno, delle
occasioni che hanno cambiato una direzione e delle esperienze che hanno trasformato un bambino incuriosito da una
pallina che rimbalzava sullo schermo in un programmatore.\newline
\textbf{C'è però un ultimo aspetto che ritengo giusto condividere con il lettore.}\newline
Questo libro è anche il risultato di un lungo lavoro di progettazione. Prima ancora di scrivere il primo capitolo ho
sentito la necessità di capire quale forma dovesse avere quest'opera, quale percorso avrebbe seguito il lettore e quali
principi avrebbero dovuto guidarne la scrittura. Per questo motivo ho trascorso molto tempo a progettare la struttura
del libro, a rivederne l'organizzazione, a mettere continuamente in discussione le scelte narrative e didattiche,
cercando un equilibrio tra il racconto di una storia personale e la costruzione di un vero percorso di
apprendimento.\newline
In questa fase mi sono confrontato a lungo con uno strumento che non sarebbe potuto esistere nell'epoca in cui questa
storia ha avuto inizio.\newline
\textbf{ChatGPT.}\newline
Non ha scritto questo libro al mio posto e non ha preso decisioni al posto mio. È stato, piuttosto, un interlocutore con
cui discutere idee, verificare la coerenza di un ragionamento, esplorare soluzioni alternative e affinare molte delle
pagine che il lettore troverà davanti a sé. Un confronto continuo, non molto diverso da quello che, in passato, avrei
potuto avere con un editor, con un collega o con un amico disposto ad ascoltare un'idea e a metterla sinceramente in
discussione.\newline
Ho scelto di raccontarlo apertamente perché anche questo fa parte della storia di questo libro.\newline
L'intelligenza artificiale non entra nella storia che sto raccontando. Quella rimane la storia di un ragazzo degli anni
Ottanta, dei suoi amici, di un insegnante e di un Commodore 64. Entra, invece, nella storia di come questo libro è
stato scritto. Sarebbe poco onesto fingere che questo dialogo non sia mai esistito, così come sarebbe stato poco
onesto, quarant'anni fa, nascondere l'importanza che ebbe il Commodore 64 nel cambiare il corso della mia vita.\newline
Ripensandoci oggi, mi colpisce una curiosa simmetria. Quando ero ragazzo utilizzavo lo strumento più moderno che potessi
immaginare: un Commodore 64. Oggi, mentre provo finalmente a mantenere quella promessa rimasta sospesa per
quarant'anni, mi accompagna uno degli strumenti più moderni che la tecnologia metta a disposizione della mia
generazione.\newline
In entrambi i casi, però, il principio rimane lo stesso. Lo strumento non sostituisce la persona, non genera la
curiosità, non prende decisioni, non costruisce un metodo e non può trasformare un ricordo in un'esperienza vissuta.
Può soltanto amplificare le possibilità di chi lo utilizza.\newline
Forse è anche questo uno dei fili invisibili che attraversano queste pagine. La tecnologia cambia, gli strumenti si
evolvono e diventano sempre più potenti, ma ciò che dà loro un significato rimane sorprendentemente immutato: la
curiosità di chi continua a porsi domande. Perché, in fondo, anche questo libro nasce proprio da una domanda.\newline
\textbf{Non so quale sarà il risultato di questo viaggio per ciascun lettore.}\newline
Forse qualcuno ritroverà emozioni che credeva dimenticate. Qualcun altro incontrerà per la prima volta una macchina
appartenente a un'altra epoca e scoprirà quanto possa ancora insegnare.\newline
Mi auguro soltanto che, arrivato all'ultima pagina, il lettore non abbia semplicemente imparato qualcosa in più sul
Commodore 64. Spero che abbia avuto la sensazione di aver partecipato davvero alla nascita di un videogioco. E,
soprattutto, che abbia scoperto il piacere di guardare un computer in modo diverso.\newline
Perché programmare, in fondo, significa proprio questo: smettere di domandarsi soltanto che cosa una macchina sappia
fare e cominciare a chiedersi come sia possibile insegnarglielo.\newline
È da domande come questa che nascono i programmatori.\newline
Ed è da una di quelle domande che, molti anni fa, è cominciato anche il mio viaggio.


\bigskip
\end{document}
